% sup2_05_S8b: Appendix on Ba3 Prop. 1, part 2, step S8, second half: the constant-gradient (zeroth-moment) form, summed internal
% positions and internal photon lines, and the scope of the identities. Input from sup2_00_intro.tex's section.

\begin{lemma}[the constant-gradient form]\label{sup2:lem:S8zero}
In \cref{sup2:lem:S8lim} let $n\ge1$ or $\beta\ne0$, and write $W^\infty=\sum_{b'}\eta^dK(b')(d\lambda)(b')$, which is possible since $W^\infty$ is linear in $d\lambda$. Then: \textup{(a)} $|K(b')|\le Ce^{-\theta\dist(b',S_0)}$, with $S_0$ the union of the supports of the $A^{(i)}$ and $X_F$; \textup{(b)} for every $c\in\R^d$, $\sum_{b'}\eta^dK(b')\,c_{\mu(b')}=0$, $\mu(b')$ the direction of $b'$. This is the zeroth-moment identity of \cite[(3.29)]{Ba3} ($c=g'(y)A'(y)$) and of \cite[(3.31)--(3.32)]{Ba3}.
\end{lemma}

\begin{proof}
(a) $K(b')$ is the sum of the connected graphs with the $\lambda$-vertex at $b'$; all other positions lie in $S_0$ and its $\eta$-neighbourhood, and in each connected graph some line joins a neighbour of $b'$ to $S_0$, so each monomial carries a factor $|C_\infty(u,u')|\le ce^{-\theta|u-u'|}$ with $u$ near $b'$ and $u'$ near $S_0$ (\cref{sup2:lem:S8CT}); the number of graphs is finite. (b) Put $\lambda(x)\defeq\sum_\mu c_\mu x_\mu$ and $\chi_R(x)\defeq\min\{1,\max\{0,2-|x|_\infty/R\}\}$. Then $\lambda_R\defeq\chi_R\lambda$ is finitely supported, $(d\lambda_R)(b')=c_{\mu(b')}$ if both ends of $b'$ lie in $\{|x|_\infty\le R\}$, and everywhere $|d\lambda_R(b')|\le|c|_\infty+2|c|_1+O(\eta/R)$, since $|\lambda|\le2R|c|_1$ on $\supp\chi_R$ and $|d\chi_R|\le1/R$. Hence $|\sum_{b'}\eta^dK(b')(c_{\mu(b')}-d\lambda_R(b'))|\le C(|c|_1)\sum_{|b'|_\infty\ge R-\eta}\eta^d|K(b')|\to0$ as $R\to\infty$ by (a), $S_0$ being fixed; and $\sum_{b'}\eta^dK(b')\,d\lambda_R(b')=0$ for every $R$ by \cref{sup2:lem:S8lim}.
\end{proof}

\Cref{sup2:lem:S8lim,sup2:lem:S8zero} are stated for finitely supported $A^{(i)}$, a finite $X_F$ and no photon integration. The uses need more: \cite[(2.28)]{Ba3} has $F=-\lambda_k\sum_{x\in\Delta}\eta^d:\!|\phi(x)|^4\!:$ summed over all cubes $\Delta$ of the class, and the degree-$0$ two-vector-leg classes of \cite[(3.31)--(3.32)]{Ba3} (``all possible connected graphs \dots\ of a fixed order'', \cite[p.~431]{Ba3}) have $\ell=2$ (\cref{sup2:lem:S7id}), so they contain internal vertices and internal photon lines at all positions. The next lemma supplies this.

\begin{lemma}[summed internal positions and internal photon lines]\label{sup2:lem:S8sum}
Work on $\xi\Z^d$ (so $\eta=\xi$) at $B=0$ with the scalar covariance $C^\xi=(-\Delta^\xi+M^2)^{-1}$ and the photon covariance $C^\xi_{\mathrm{ph}}=(-\Delta^\xi+M^2)^{-1}\otimes\mathbf 1_d$, $M^2=1$ (\cite[p.~442]{Ba3}: ``replace the propagators $G_{j_0}(0)$, $G_{j_0}$ by $C^\xi$''). Fix finitely many external vector legs at bonds of a finite set $S_0$, gauge-invariant vertex densities $f_r$ (polynomials in $\phi(x)\cdot P(q)\phi(x)$, for instance $-\lambda_k:\!|\phi|^4\!:$ and $-\frac12\delta m^2:\!|\phi|^2\!:$), and an order in $e_k$ and in the $f_r$. Let $W^{\mathrm{sum}}$ be the sum, over all connected graphs of that order with full lines $C^\xi$ (scalar) and $C^\xi_{\mathrm{ph}}$ (photon), of the graphs with one $\lambda$-vertex (linear in $d\lambda$), the external legs at $S_0$, the $f_r$-vertices summed over all $x\in\xi\Z^d$ with weight $\xi^d$, and the $A$-vertices of $H(A)$ (the vertices \cite[(1.8), (1.10)]{Ba3} at $B=0$) summed over all bonds, their $A$-legs either external (at $S_0$) or contracted in pairs by $C^\xi_{\mathrm{ph}}$. Then: \textup{(a)} for finitely supported $\lambda$, $W^{\mathrm{sum}}=0$; \textup{(b)} if $S_0\ne\emptyset$, then writing $W^{\mathrm{sum}}=\sum_{b'}\xi^dK^{\mathrm{sum}}(b')(d\lambda)(b')$ one has $|K^{\mathrm{sum}}(b')|\le Ce^{-\theta'\dist(b',S_0)}$ for some $\theta'>0$, and $\sum_{b'}\xi^dK^{\mathrm{sum}}(b')c_{\mu(b')}=0$ for every $c\in\R^d$.
\end{lemma}

With $S_0=\emptyset$ the class is translation invariant, $K^{\mathrm{sum}}(b')$ depends only on the direction of $b'$, and (b) is not meaningful as written; (a) needs no such hypothesis. The uses on \cite[p.~442]{Ba3}, \cite[(3.29), (3.31)--(3.32)]{Ba3}, have external vector legs.

\begin{proof}
\emph{Step 1: the polynomial identity, pointwise in the photon field.} On a torus $T_\ell$ ($\ell\in\xi\N$) let $A$ be any real bond field. For real $e_k$ near $0$ the functions $e_k\mapsto Z_\beta(A;e_k)$ of \cref{sup2:lem:S8torus} are smooth ($H(A)\ge M^2$), and their Taylor coefficients $P_{r,\beta}(A)\defeq\frac1{r!}\partial^r_{e_k}Z_\beta(A;e_k)|_{e_k=0}$ are polynomials in the finitely many $A_b$ (the expansion of $U(A)=\exp(e_k\xi qA)$ gives vertices with $A$-legs, each $e_k^r$ coming with $A^r$). The invariance $Z_\beta(A+d\lambda;e_k)=Z_\beta(A;e_k)$ for all real $e_k$ near $0$ and all $A$, $\lambda$ (\cref{sup2:lem:S8torus}, $T_\lambda$ depending on $e_k$) gives $P_{r,\beta}(A+d\lambda)=P_{r,\beta}(A)$ as polynomial identities, for every $r$, $\beta$, $A$ and $\lambda$ on $T_\ell$.
\emph{Step 2: integrating the photon.} Put $A=A_{\mathrm{ext}}+A_{\mathrm{int}}$ with $A_{\mathrm{int}}$ Gaussian of covariance $C^\ell_{\mathrm{ph}}$, the torus photon covariance at $M^2=1$ (a nondegenerate finite-dimensional Gaussian). The $A_{\mathrm{int}}$-expectation of the identity of Step 1 (a polynomial in $A_{\mathrm{int}}$, so a finite Wick sum) gives $\E[P_{r,\beta}(A_{\mathrm{ext}}+d\lambda+A_{\mathrm{int}})]=\E[P_{r,\beta}(A_{\mathrm{ext}}+A_{\mathrm{int}})]$, and by Wick's theorem $\E[P_{r,\beta}(\cdot)]$ is the sum over graphs with internal photon lines $C^\ell_{\mathrm{ph}}$. Every coefficient of the formal logarithm in $(e_k,\tau)$ is a polynomial in the $\E[P_{r,\beta}]/\E[P_{0,0}]$ ($P_{0,0}>0$ is $A$-independent), so it is gauge invariant in $A_{\mathrm{ext}}$ and its mixed $A_{\mathrm{ext}}$-$d\lambda$ derivatives vanish. These derivatives are the connected graphs of the class with scalar lines $C^\ell$, photon lines $C^\ell_{\mathrm{ph}}$ and all internal positions summed over $T_\ell$ (the $f_r$-vertices entering through $F_r=\sum_x\xi^df_r(\phi(x))$). This is (a) on $T_\ell$, with $\lambda$ finitely supported and lifted as in \cref{sup2:lem:S8lim}.
\emph{Step 3: absolute convergence, uniformly in $\ell$.} By \cref{sup2:lem:S8CT}, $|C^\xi(u,u')|,|C^\xi_{\mathrm{ph}}(u,u')|\le ce^{-\theta|u-u'|}$ and $|C^\ell(u,u')|,|C^\ell_{\mathrm{ph}}(u,u')|\le ce^{-\theta d_\ell(u,u')}$ uniformly in $\ell$, $d_\ell$ the torus distance. Every graph of the class is connected and has a fixed vertex (an external leg at $S_0$, or the $\lambda$-vertex on the finite $\supp d\lambda$). Choose a spanning tree rooted at a fixed vertex and bound the non-tree lines by $c$; changing variables to the tree differences $z_i$ (minimal representatives on $T_\ell$) bounds each graph by $\prod_ice^{-\theta|z_i|}$ times constants, summable over $(\xi\Z^d)^{\#\text{tree lines}}$ uniformly in $\ell$.
\emph{Step 4: the limit $\ell\to\infty$.} For fixed $z$'s all positions stay within a fixed distance of the root, so for $\ell$ large the torus distances equal the lattice distances and $C^\ell\to C^\xi$, $C^\ell_{\mathrm{ph}}\to C^\xi_{\mathrm{ph}}$ there (\cref{sup2:lem:S8lim}). By dominated convergence (Step 3) each graph converges to its value on $\xi\Z^d$, with absolutely convergent position sums; the number of graphs of the fixed order is finite. So $W^{\mathrm{sum}}=\lim_\ell W^{\mathrm{sum},T_\ell}=0$, which is (a).
\emph{Step 5: (b).} As in \cref{sup2:lem:S8zero}(a): in each connected graph the $\lambda$-vertex at $b'$ is joined to $S_0$ by a path of lines; summing the intermediate positions along this path with $\sum_ze^{-\theta|b'-z|}e^{-\theta|z-s|}\le C_{\theta'}e^{-\theta'|b'-s|}$ ($\theta'<\theta$) and the other positions by Step 3 gives the decay of $K^{\mathrm{sum}}$. The cut-off $\lambda_R=\chi_R\lambda$ of \cref{sup2:lem:S8zero}(b) and dominated convergence then give the zeroth-moment identity exactly as there.
\end{proof}

\begin{remark}[scope of the identities]\label{sup2:rem:S8scope}
\textup{(a)} \emph{Why the limit is needed.} On $T_\ell$ a linear $\lambda$ is not a function on the torus, and a constant $d\lambda=c$ is a flat connection with holonomy, not gauge-trivial; so \cref{sup2:lem:S8zero}(b) is false on a torus in general, the defect being a sum over images winding around the torus, exponentially small in $\ell M$. Ba{\l}aban's ``by taking a limit'' is essential, and it is exactly \cref{sup2:lem:S8lim,sup2:lem:S8zero,sup2:lem:S8sum}: the torus identity with finitely supported $\lambda$, the limit, then the extension to linear $\lambda$ on $\xi\Z^d$.
\textup{(b)} \emph{Full lines only.} \Cref{sup2:lem:S8torus,sup2:lem:S8lim,sup2:lem:S8zero,sup2:lem:S8sum} are identities between sums over all graphs of a given order with every line the full propagator. They say nothing about a single scale assignment or a fixed ordering of slices; after the decomposition \cite[(2.6)]{Ba3} the per-ordering, per-slice analogue is false. \cite{Ba3} applies \cite[(2.26)]{Ba3} only after ``we resum over proper orderings and $j$-indices, we replace the propagators $G_{j_0}(0)$, $G_{j_0}$ by $C^\xi$'' \cite[p.~442]{Ba3}, which is the full-line form; that this resummation and replacement is legitimate at the places of \cite[Sect.~3]{Ba3} is part of the renormalisation step (Step~3(D$_4$) of \cref{app:a06_BALtorus}), not of S8.
\textup{(c)} \emph{Background and averaging terms.} \Cref{sup2:lem:S8cov,sup2:lem:S8torus,sup2:lem:S8lim,sup2:lem:S8zero} hold at any background $B$; \cite{Ba3} uses $B=0$, after the gauge transformation removing $B_0$. If $H(A)$ contains a covariant averaging term $a_kP_k(B+A)$ (the propagators $G_k$), \cref{sup2:lem:S8cov} still holds, since $Q_k(A+d\lambda)T_\lambda^{-1}=T^{\mathrm{unit}}_\lambda Q_k(A)$ by transport along the contours; so the identities hold with the vertices \cite[(1.13)--(1.15)]{Ba3} added. The printed and the consumed identities use $C_{M^2}$, which needs no such remark. \textup{(d)} The constants of \cref{sup2:lem:S8sum} depend on $\xi$ and the order; only the limits $\ell\to\infty$ and $R\to\infty$ are taken, at fixed $\xi$, which is how the identities are used.
\end{remark}
